\documentclass[12pt]{article}
\usepackage[a4paper,margin=1in]{geometry}
\usepackage[utf8]{inputenc}
\usepackage[T1]{fontenc}
\usepackage{amsmath,amssymb}
\usepackage{textcomp}
\usepackage{parskip}
\begin{document}
\begin{center}
{\Large \textbf{2003 STEP II}}
\end{center}

\subsection*{Page 1}
\begin{flushleft}
STEP II, 2003 2\par
\bigskip\textbf{Section A: Pure Mathematics}\par
\bigskip\textbf{Question 1.} Consider the equations\par
\[ax y z = 3 ,\]
- -\par
\[2ax y 3z = 7 ,\]
- -\par
\[3ax y 5z = b ,\]
- -\par
where a and b are given constants.\par
\textbf{(i)} In the case a = 0 , show that the equations have a solution if and only if b = 11 .\par
\textbf{(ii)} In the case a = 0 and b = 11 show that the equations have a solution with z = λ for\par
any given number λ .\par
\textbf{(iii)} In the case a = 2 and b = 11 find the solution for which x2 + y2 + z2 is least.\par
\textbf{(iv)} Find a value for a for which there is a solution such that x > 106 and y2 + z2 < 1 .\par
\[2 Write down a value of θ in the interval \pi/4 < θ < \pi/2 that satisfies the equation\]
\[4 cos θ + 2\sqrt\{3\} sin θ = 5 .\]
Hence, or otherwise, show that\par
\[\pi = 3 arccos(5/\sqrt\{28\}) + 3 arctan(\sqrt\{3\}/2) .\]
Show that\par
\[\pi = 4 arcsin(7\sqrt\{2\}/10) 4 arctan(3/4) .\]
-\par
\bigskip\textbf{Question 3.} Prove that the cube root of any irrational number is an irrational number.\par
\[Let u = 5\]
1/(3n)\par
\[. Given that \sqrt\{3\} 5 is an irrational number, prove by induction that u is an\]
n n\par
irrational number for every positive integer n.\par
Hence, or otherwise, give an example of an infinite sequence of irrational numbers which\par
converges to a given integer m .\par
[An irrational number is a number that cannot be expressed as the ratio of two integers.]\par
\end{flushleft}
\subsection*{Page 2}
\begin{flushleft}
STEP II, 2003 3\par
\bigskip\textbf{Question 4.} The line y = d , where d > 0 , intersects the circle x2 + y2 = R2 at G and H. Show that the\par
area of the minor segment GH is equal to\par
d\par
\[R^2 arccos d R^2 d^2 . ( )\]
\lambda R \mu - - ∗\par
p\par
In the following cases, the given line intersects the given circle. Determine how, in each case,\par
the expression ( ) should be modified to give the area of the minor segment.\par
∗\par
\textbf{(i)} Line: y = c ; circle: (x a)2 + (y b)2 = R2 .\par
- -\par
\textbf{(ii)} Line: y = mx + c ; circle: x2 + y2 = R2 .\par
\textbf{(iii)} Line: y = mx + c ; circle: (x a)2 + (y b)2 = R2 .\par
- -\par
\bigskip\textbf{Question 5.} The position vectors of the points A , B and P with respect to an origin O are ai , bj and\par
li + mj + nk , respectively, where a, b, and n are all non-zero. The points E, F , G and H\par
are the midpoints of OA, BP , OB and AP , respectively. Show that the lines EF and GH\par
intersect.\par
Let D be the point with position vector dk, where d is non-zero, and let S be the point of\par
intersection of EF and GH. The point T is such that the mid-point of DT is S. Find the\par
position vector of T and hence find d in terms of n if T lies in the plane OAB.\par
\bigskip\textbf{Question 6.} The function f is defined by\par
\[f (x) = x 1 ,\]
| - |\par
where the domain is R , the set of all real numbers. The function g = f n, with domain R , so\par
n\par
for example g (x) = f (f (f (x))) . In separate diagrams, sketch graphs of g , g , g and g .\par
3 1 2 3 4\par
The function h is defined by\par
\[\pix\]
\[h (x) = sin ,\]
\{ 2 \{\par
\{ \{\par
where the domain is R . Show that if n is eve\{n, \{\par
n 2n n\par
\[h (x) g (x) dx = .\]
\[Z - n \pi - 2\]
\end{flushleft}
\subsection*{Page 3}
\begin{flushleft}
STEP II, 2003 4\par
\bigskip\textbf{Question 7.} Show that, if n > 0 , then\par
\infty\par
ln x 2\par
\[dx = .\]
\[Z e^1/n xn+1 n2e\]
ln x\par
You may assume that 0 as x .\par
x \to \to \infty\par
\[Explain why, if 1 < a < b , then\]
\infty \infty\par
ln x ln x\par
\[dx < dx .\]
Z xn+1 Z xn+1\par
b a\par
Deduce that\par
N\par
1 e\par
\infty\par
1 x\par
-N\par
\[< - ln x dx ,\]
\[nX=1 n^2 2 Z e^1/N \lambda x^2 - x \mu\]
where N is any integer greater than 1.\par
8 It is given that y satisfies\par
\[dy t^2 3t + 2\]
\[+ k - y = 0 ,\]
dt \lambda t + 1 \mu\par
where k is a constant, and y = A when t = 0 , where A is a positive constant. Find y in terms\par
of t , k and A .\par
\[Show that y has two stationary values whose ratio is (3/2)6ke -5k/2.\]
\[Describe the behaviour of y as t + for the case where k > 0 and for the case where k < 0 .\]
\to \infty\par
\[In separate diagrams, sketch the graph of y for t > 0 for each of these cases.\]
\end{flushleft}
\subsection*{Page 4}
\begin{flushleft}
STEP II, 2003 5\par
\bigskip\textbf{Section B: Mechanics}\par
\[9 AB is a uniform rod of weight W . The point C on AB is such that AC > CB . The rod is\]
in contact with a rough horizontal floor at A and with a cylinder at C . The cylinder is fixed\par
to the floor with its axis horizontal. The rod makes an angle α with the horizontal and lies in\par
a vertical plane perpendicular to the axis of the cylinder. The coefficient of friction between\par
the rod and the floor is tan λ and the coefficient of friction between the rod and the cylinder\par
is tan λ .\par
Show that if friction is limiting both at A and at C, and α = λ λ , then the frictional force\par
2 1\par
6 -\par
acting on the rod at A has magnitude\par
W sin λ sin(α λ )\par
1 2\par
- .\par
sin(α + λ λ )\par
1 2\par
-\par
\bigskip\textbf{Question 10.} A bead B of mass m can slide along a rough horizontal wire. A light inextensible string of\par
length 2` has one end attached to a fixed point A of the wire and the other to B . A particle\par
P of mass 3m is attached to the mid-point of the string and B is held at a distance ` from A .\par
The bead is released from rest.\par
Let a and a be the magnitudes of the horizontal and vertical components of the intial\par
1 2\par
acceleration of P . Show by considering the motion of P relative to A , or otherwise, that\par
\[a = \sqrt\{3a\} . Show also that the magnitude of the intial acceleration of B is 2a .\]
1 2 1\par
Given that the frictional force opposing the motion of B is equal to (\sqrt3/6)R, where R is the\par
normal reaction between B and the wire, show that the magnitude of the intial acceleration\par
of P is g/18 .\par
\textasciicircum{}\circ\par
\bigskip\textbf{Question 11.} A particle P is projected with speed V at an angle of elevation α (> 45 ) , from a point in\par
a horizontal plane. Find T , the flight time of P , in terms of α, V and g . Show that the time\par
1 1\par
\textasciicircum{}\circ\par
after projection at which the direction of motion of P first makes an angle of 45 with the\par
horizontal is 1 (1 cot α)T .\par
2 - 1\par
A particle P is projected under the same conditions. When the direction of the motion of P\par
2 2\par
\textasciicircum{}\circ\par
first makes an angle of 45 with the horizontal, the speed of P is instantaneously doubled. If\par
T is the total flight time of P , show that\par
2 2\par
2T\par
\[2 = 1 + cot α + 1 + 3 cot2 α .\]
T\par
1 p\par
\end{flushleft}
\subsection*{Page 5}
\begin{flushleft}
STEP II, 2003 6\par
SECTION C: Probability and Statistics\par
\bigskip\textbf{Question 12.} The life of a certain species of elementary particles can be described as follows. Each particle\par
has a life time of T seconds, after which it disintegrates into X particles of the same species,\par
where X is a random variable with binomial distribution B(2, p) . A population of these\par
particles starts with the creation of a single such particle at t = 0 . Let X be the number of\par
n\par
\[particles in existence in the time interval nT < t < (n + 1)T , where n = 1 , 2 , . . ..\]
\[Show that P(X = 2 and X = 2) = 6p4q2 , where q = 1 p . Find the possible values of p if\]
1 2\par
-\par
\[it is known that P(X = 2 X = 2) = 9/25 .\]
1 2\par
|\par
\[Explain briefly why E(X n ) = 2pE(X n-1 ) and hence determine E(X n ) in terms of p. Show\]
that for one of the values of p found above lim n\to\infty E(X n ) = 0 and that for the other\par
\[lim n\to\infty E(X n ) = + .\]
\infty\par
\bigskip\textbf{Question 13.} The random variable X takes the values k = 1, 2, 3, . . . , and has probability distribution\par
λke -λ\par
\[P(X = k) = A ,\]
k!\par
\[where λ is a positive constant. Show that A = (1 e\]
-λ) -1\par
. Find the mean µ in terms of λ\par
-\par
and show that\par
\[Var (X) = µ(1 µ + λ) .\]
-\par
\[Deduce that λ < µ < 1 + λ .\]
Use a normal approximation to find the value of P (X = λ) in the case where λ = 100 , giving\par
your answer to 2 decimal places.\par
\bigskip\textbf{Question 14.} The probability of throwing a 6 with a biased die is p . It is known that p is equal to one or\par
\[other of the numbers A and B where 0 < A < B < 1 . Accordingly the following statistical\]
test of the hypothesis H : p = B against the alternative hypothesis H : p = A is performed.\par
0 1\par
The die is thrown repeatedly until a 6 is obtained. Then if X is the total number of throws,\par
H is accepted if X 6 M , where M is a given positive integer; otherwise H is accepted. Let\par
0 1\par
α be the probability that H is accepted if H is true, and let β be the probability that H is\par
1 0 0\par
accepted if H is true.\par
\[Show that β = 1 αK , where K is independent of M and is to be determined in terms of A\]
-\par
and B . Sketch the graph of β against α .\par
\end{flushleft}
\end{document}
